\documentclass[12pt]{article}
\usepackage[utf8]{inputenc}
\usepackage[T1]{fontenc}
\usepackage{hyperref}
\title{Chut Language Resources for K-12 Curriculum Development}
\author{Compiled by Claude Sonnet 5.5 (claude-sonnet-5-5) for Andr\'as Kornai}
\date{}
\begin{document}
\maketitle

\section{Introduction}
Chut (ISO 639-3 scb) is a Vietic (Austroasiatic) language cluster, also called Rục-Sách. It is spoken in Quảng Bình, Vietnam, and in Khammouane Province, Laos. Wikipedia gives about 1,300 speakers (2007). ``Chứt'' is the official Vietnamese ethnic label for five groups: Arem, Mã Liêng, Mày, Rục and Sách. The Arem speak a separate language. The 2019 census counted 7,513 Chứt people. This is a search-based survey. Very little exists for this language, and no school materials were found.

\section{Primary Educational Resources}
No primers, textbooks or school programmes in Chut were found in the searches. Nothing indicates that the language is taught in schools.

\section{Linguistic and Reference Materials}
\begin{itemize}
\item Overview of the Chut languages (Wikipedia): \url{https://en.wikipedia.org/wiki/Chut_languages}
\item Chứt people (Wikipedia): \url{https://en.wikipedia.org/wiki/Ch%E1%BB%A9t_people}
\item May people (Wikipedia). The Wikipedia references cite the Babaev and Samarina (2021, Brill) grammar of May: \url{https://en.wikipedia.org/wiki/May_people}
\item Numerals data for the Sach dialect (MPI-EVA): \url{https://lingweb.eva.mpg.de/channumerals/Chut.htm}
\end{itemize}
No dictionary of Sách or Rục was found.

\section{Religious and Community Resources}
\begin{itemize}
\item Article on the Chứt ethnic group, French edition (Nhân Dân): \url{https://special.nhandan.vn/dt_chut_fr/index.html}
\item Page on the Chut people (vietnam.vn): \url{https://www.vietnam.vn/nguoi-chut}
\end{itemize}
No scripture translations and no community language organisations were found.

\section{Digital and Audio Resources}
\begin{itemize}
\item Sach audio corpus from the Pangloss Collection (LACITO-CNRS), recorded in Yên Hợp hamlet, Thượng Hóa commune, Quảng Bình. It is listed on the Endangered Languages Project: \url{https://endangeredlanguages.com/lang/1215/samples/10374}
\item Arem entry (Verbix Endangered Languages): \url{https://docs.verbix.com/EndangeredLanguages/Arem}
\end{itemize}

\section{Curriculum Development Context}
Chut has no known standard orthography and no known school use. A curriculum would have to start from documentation work: building on the existing recordings, agreeing an orthography with the communities, and compiling a lexicon. The varieties (Rục, Sách, Mày, Maleng) differ, and Arem is a separate language, so each would need its own treatment. Vietnamese is the likely language of instruction.

\end{document}